\documentclass[11pt,a4paper]{article}

\usepackage[utf8]{inputenc}
\usepackage[T1]{fontenc}
\usepackage{lmodern}
\usepackage[margin=2.5cm]{geometry}
\usepackage{amsmath,amssymb,mathtools}
\usepackage{booktabs}
\usepackage{xcolor}
\usepackage{pgfplots}
\usepackage{pgfplotstable}
\pgfplotsset{compat=1.17}
\usepgfplotslibrary{groupplots}
\usepackage[colorlinks=true,linkcolor=blue!50!black,urlcolor=blue!50!black]{hyperref}
\usepackage[capitalise,noabbrev]{cleveref}
\usepackage{url}
\DeclareUrlCommand\code{\urlstyle{tt}}
\emergencystretch=2em

\newcommand{\dd}{\,\mathrm{d}}
\newcommand{\pos}[1]{\langle #1\rangle_{+}}


% slope triangle: \slopetriangle{x0}{y0}{x1}{slope}{label}
\newcommand{\slopetriangle}[5]{%
  \draw[gray] (axis cs:#1,#2) -- (axis cs:#3,#2) -- (axis cs:#3,{#2*(#3/#1)^(#4)}) -- cycle;
  \node[font=\scriptsize, gray, anchor=west] at (axis cs:#3,{#2*(#3/#1)^(#4/2)}) {#5};
}

\title{Frictionless contact with rigid obstacles in ExaDG:\\
penalty formulation and validation}
\author{Branch \code{structure-contact}}
\date{\today}

\begin{document}
\maketitle

\begin{abstract}
  We summarise the frictionless penalty contact of a hyperelastic body with rigid obstacles that
  was added to the structure module of ExaDG, and its validation. A contact patch test is
  reproduced to round-off in two and three dimensions, for several polynomial degrees and all
  available preconditioners. For the Hertz problem, the discretisation and penalty errors
  converge, and the remaining deviation from the Hertz solution scales like the square root of
  the indentation, i.e.\ like the strain: it is the expected difference between the finite-strain
  model and the linear Hertz theory. Newton's method converges quadratically once the active set
  is settled.
\end{abstract}

\section{Formulation}

Let $\Gamma_C$ be the part of the boundary of the reference configuration $\Omega_0$ that may come
into contact with a rigid obstacle. The obstacle is given by a signed distance function $\phi$:
for a material point $X\in\Gamma_C$ with current position $x = X + d(X)$,
\begin{equation}
  g = \phi(x), \qquad n = \nabla\phi(x), \qquad \kappa = \nabla\nabla\phi(x)
\end{equation}
are the gap ($g<0$: penetration), the unit normal pointing from the obstacle towards the body and
the curvature of the obstacle. The penalty contact pressure and the contact term of the weak form
(measured per reference area) are
\begin{equation}
  \lambda = \varepsilon \pos{-g}, \qquad
  r_C(d; v) = -\int_{\Gamma_C} \lambda\, n\cdot v \dd A ,
\end{equation}
and the consistent linearisation in direction $\delta d$ reads
\begin{equation}
  \delta r_C = \int_{\Gamma_C} \Bigl(\varepsilon\, H(-g)\,(n\cdot\delta d)\, n
  - \lambda\, \kappa\, \delta d\Bigr)\cdot v \dd A ,
\end{equation}
with the Heaviside function $H$ (points with $g=0$ are counted as active). The term is symmetric,
so the conjugate gradient method can still be used for the linearised problems.

\paragraph{Implementation.}
The obstacle is prescribed per boundary id via \code{BoundaryDescriptor::contact_bc} (a
\code{dealii::Function} providing value, gradient and Hessian, evaluated at the operator time so
that obstacles may move), the penalty via \code{Parameters::contact_penalty_parameter}. The
residual is added in the boundary face loop of \code{NonLinearOperator}; the linearisation is
hooked into \code{reinit\_boundary\_face\_derived}/\code{do\_boundary\_integral\_continuous}, so
the matrix-free operator, its diagonal, the matrix-based assembly (AMG) and all multigrid levels
contain it. Contact force, active area, pressure and gap are available in the postprocessor. All
results below are computed with the application \code{structure/contact\_rigid\_obstacle},
St.\,Venant--Kirchhoff material with $E=1$, $\nu=0.3$ (plane strain in 2d), quasi-static load
stepping and quadratic elements unless stated otherwise.

\section{Contact patch test}

A block $[0,2]\times[0,1]$ ($\times[0,1]$ in 3d) rests on the rigid plane $y=0$ and is loaded by a
uniform traction $p$ on the top face; symmetry conditions on $x=0$ (and $z=0$) remove the
tangential rigid body modes. The exact solution of the discrete penalty problem is a homogeneous
deformation with uniform contact pressure $\lambda=p$, penetration $g=-p/\varepsilon$, contact
force $p\,|\Gamma_C|$ and active area $|\Gamma_C| = 2$, independently of mesh, degree and strain
magnitude. \Cref{tab:patch} reports the largest relative error of these four quantities
($\varepsilon=10^3$): all of them agree with the exact values in all ten printed digits, including
$p=0.15$, which corresponds to a compressive strain of about $15\,\%$.

\begin{table}[ht]
  \centering
  \pgfplotstableread{data/patch.dat}\patchtab
  \pgfplotstabletypeset[
    columns={name,dim,p,dofs,max_rel_err,newton},
    columns/name/.style={string type, column name={case},
      string replace={D_2d_q2_amg}{2d, $Q_2$, AMG},
      string replace={D_2d_q3_mg}{2d, $Q_3$, matrix-free multigrid},
      string replace={D_2d_q2_mb}{2d, $Q_2$, matrix-based multigrid},
      string replace={D_2d_q2_large}{2d, $Q_2$, multigrid, large strain},
      string replace={D_3d_q2_mg}{3d, $Q_2$, matrix-free multigrid},
      string replace={D_3d_q3_large}{3d, $Q_3$, multigrid, large strain},
      column type=l},
    columns/dim/.style={column name={$d$}},
    columns/p/.style={column name={$p$}, fixed, precision=2},
    columns/dofs/.style={column name={DoFs}},
    columns/max_rel_err/.style={column name={max.\ rel.\ error}, fixed, precision=0},
    columns/newton/.style={column name={Newton it.\ (total)}},
    every head row/.style={before row=\toprule, after row=\midrule},
    every last row/.style={after row=\bottomrule},
  ]{\patchtab}
  \caption{Contact patch test: largest relative error of contact force, active area, pressure and
  gap with respect to the exact values (0: identical in all ten printed digits).}
  \label{tab:patch}
\end{table}

\section{Hertz problem}

A rigid cylinder (2d, plane strain) of radius $R=1$ is pressed by an indentation $\delta$ into the
elastic body $[0,4]\times[-4,0]$; only one half is computed (symmetry on $x=0$, clamped bottom).
The mesh is refined locally towards the initial contact point; $h$ denotes the element size in the
contact zone. For the computed load $P$ (per unit length, full cylinder), the Hertz solution is
\begin{equation}
  a = \sqrt{\frac{4PR}{\pi E^*}}, \qquad p_0 = \frac{2P}{\pi a}, \qquad
  p(r) = p_0\sqrt{1-r^2/a^2}, \qquad E^* = \frac{E}{1-\nu^2}.
\end{equation}
Three error sources are present: the discretisation ($h$), the penalty ($\varepsilon$), and the
modelling difference between the finite-strain hyperelastic problem and the linear Hertz theory.
They are separated by the three studies in \cref{fig:hertz}.

\begin{figure}[p]
  \centering
  \begin{tikzpicture}
    \begin{axis}[width=0.48\textwidth, height=0.42\textwidth, xmode=log, ymode=log, grid=major,
        grid style={gray!25}, legend style={font=\footnotesize, fill=white, fill opacity=0.85,
        text opacity=1}, tick label style={font=\footnotesize}, label style={font=\small},
        xlabel={indentation $\delta/R$}, ylabel={relative error},
        title={\small (a) deviation from Hertz, $\varepsilon=100/\delta$},
        legend pos=south east, ymin=1e-3, ymax=2e-1]
      \addplot+[mark=*] table[x=delta, y=err_p0] {data/indentation.dat};
      \addlegendentry{$p(0)$}
      \addplot+[mark=triangle*] table[x=delta, y=err_l2_int] {data/indentation.dat};
      \addlegendentry{$p(r)$, $r<0.8a$}
      \addplot+[mark=square*] table[x=delta, y=err_a] {data/indentation.dat};
      \addlegendentry{$a$}
      \slopetriangle{3e-4}{2.5e-3}{3e-3}{0.5}{$1/2$}
    \end{axis}
  \end{tikzpicture}\hfill
  \begin{tikzpicture}
    \begin{axis}[width=0.48\textwidth, height=0.42\textwidth, xmode=log, ymode=log, grid=major,
        grid style={gray!25}, legend style={font=\footnotesize, fill=white, fill opacity=0.85,
        text opacity=1}, tick label style={font=\footnotesize}, label style={font=\small},
        xlabel={$h/a$}, ylabel={relative error},
        title={\small (b) mesh, $\delta=10^{-3}$, $\varepsilon=10^5$},
        legend pos=north west, ymin=5e-6, ymax=1e2]
      \addplot+[mark=*] table[x=h_over_a, y=err_P] {data/mesh.dat};
      \addlegendentry{$P$}
      \addplot+[mark=triangle*] table[x=h_over_a, y=err_profile_int] {data/mesh.dat};
      \addlegendentry{$p(r)$, $r<0.8a$}
      \addplot+[mark=square*] table[x=h_over_a, y=err_profile] {data/mesh.dat};
      \addlegendentry{$p(r)$, $r<2a$}
      \addplot+[mark=diamond*, dashed, black] table[x=h_over_a, y=err_hertz_int] {data/mesh.dat};
      \addlegendentry{$p(r)$, $r<0.8a$, vs.\ Hertz}
    \end{axis}
  \end{tikzpicture}

  \vspace{1em}
  \begin{tikzpicture}
    \begin{axis}[width=0.48\textwidth, height=0.42\textwidth, xmode=log, ymode=log, grid=major,
        grid style={gray!25}, legend style={font=\footnotesize, fill=white, fill opacity=0.85,
        text opacity=1}, tick label style={font=\footnotesize}, label style={font=\small},
        xlabel={penalty $\varepsilon$}, ylabel={relative error / penetration},
        title={\small (c) penalty, $\delta=10^{-3}$, $h/a\approx 0.05$},
        legend pos=north east, ymax=1e2]
      \addplot+[mark=*] table[x=eps, y=err_P] {data/penalty.dat};
      \addlegendentry{$P$ vs.\ $\varepsilon=10^7$}
      \addplot+[mark=triangle*] table[x=eps, y=penetration_over_delta] {data/penalty.dat};
      \addlegendentry{$\max|g|/\delta$}
      \addplot+[mark=diamond*, dashed, black] table[x=eps, y=err_hertz_int] {data/penalty.dat};
      \addlegendentry{$p(r)$, $r<0.8a$, vs.\ Hertz}
      \slopetriangle{1e5}{1e-5}{1e6}{-1}{$-1$}
    \end{axis}
  \end{tikzpicture}\hfill
  \begin{tikzpicture}
    \begin{axis}[width=0.48\textwidth, height=0.42\textwidth, grid=major, grid style={gray!25},
        xlabel={$r/a$}, ylabel={$p/p_0$}, xmin=0, xmax=1.5, ymin=-0.05, ymax=1.15,
        title={\small (d) pressure profile, $\delta=10^{-3}$},
        legend style={font=\footnotesize}, legend pos=north east,
        tick label style={font=\footnotesize}, label style={font=\small}]
      \addplot[black, thick, domain=0:1, samples=200] {sqrt(max(1-x^2,0))};
      \addlegendentry{Hertz}
      \addplot[only marks, mark=o, mark size=1.2pt, blue] table[x=r_over_a, y=p_over_p0]
        {data/profile_L6.dat};
      \addlegendentry{$h/a\approx 0.4$}
      \addplot[only marks, mark=*, mark size=0.6pt, red] table[x=r_over_a, y=p_over_p0]
        {data/profile_L11.dat};
      \addlegendentry{$h/a\approx 0.01$}
    \end{axis}
  \end{tikzpicture}
  \caption{Hertz problem (2d). (a) Error with respect to the Hertz solution for decreasing
  indentation ($a/h$ between 10 and 30): it decays like $\delta^{1/2}$, i.e.\ like the strain
  $p_0/E^*$. (b) Mesh convergence of load and pressure profile with respect to the finest mesh
  ($h/a = 0.0135$); the error with respect to Hertz (dashed) saturates at the modelling error
  of (a). (c) Penalty convergence with respect to $\varepsilon=10^7$: penetration and load error
  decay like $\varepsilon^{-1}$. (d) Contact pressure at the quadrature points of the contact
  boundary.}
  \label{fig:hertz}
\end{figure}

\paragraph{Indentation (\cref{fig:hertz}a).}
With $a/h$ between 10 and 30 and a penalty $\varepsilon=100/\delta$ that keeps the penetration
below $10^{-3}\delta$, the error of the pressure in the centre decreases with slope $1/2$ in
$\delta$, from $7.6\,\%$ at $\delta=3\cdot10^{-2}$ to $0.36\,\%$ at $\delta=10^{-4}$. Since
$p_0/E^* \propto \sqrt{\delta/R}$, this is a first-order error in the strain, i.e.\ the expected
discrepancy between the geometrically nonlinear St.\,Venant--Kirchhoff model and the linear Hertz
theory, and not a discretisation error, as confirmed by (b) and (c). The interior pressure profile
follows the same trend down to $\delta = 10^{-3}$; for the two smallest indentations the mesh has
only $a/h\approx 10$ and the profile error stays at the discretisation level of about $1\,\%$. The
contact radius is measured through the active area, which is resolved only up to the spacing of
the quadrature points, and is therefore less regular ($0.9$--$6\,\%$).

\paragraph{Mesh (\cref{fig:hertz}b).}
The load converges from $1.6\cdot10^{-4}$ to $2\cdot10^{-6}$ with respect to the finest mesh. The
pressure in the interior of the contact zone ($r<0.8a$) is pre-asymptotic as long as the contact
zone is covered by fewer than about five elements ($h/a>0.2$), and then converges rapidly, from
$10^{-3}$ at $h/a=0.11$ to $1.5\cdot10^{-5}$ at $h/a=0.027$. The error of the whole profile
($r<2a$) is dominated by the elements cut by the boundary of the contact zone, where the pressure
has a square-root behaviour and the active set is resolved only at quadrature points; it
decreases irregularly to $5\cdot10^{-3}$. The interior error with respect to Hertz levels off at
$0.87\,\%$, the modelling error for $\delta=10^{-3}$ (consistent with (a)).

\paragraph{Penalty (\cref{fig:hertz}c).}
The maximal penetration equals $p_0/\varepsilon$ and the load converges like $\varepsilon^{-1}$
to the one with the largest penalty, as expected for the penalty method. The interior pressure
reaches the modelling floor for $\varepsilon\gtrsim10^5$ ($\varepsilon h/E\approx 10^2$). For
$\varepsilon=10^7$ ($\varepsilon h/E\approx 10^4$) the error increases again slightly
($0.86\,\%\to0.93\,\%$) and the pressure starts to oscillate between neighbouring quadrature points:
for $\varepsilon\to\infty$ the penalty enforces $g=0$ at every quadrature point of the face, which
are more than the displacement degrees of freedom, so the constraint becomes over-determined
(locking). The penalty must therefore be chosen in a window, here
$10^4\lesssim\varepsilon\lesssim10^6$.

\paragraph{Three dimensions.}
For a rigid sphere ($\delta=10^{-3}$, quarter domain, matrix-free multigrid), the computed load
agrees with the Hertz value $P=\tfrac43 E^*\sqrt{R}\,\delta^{3/2}$ within $0.2\,\%$.
\Cref{tab:hertz3d} shows that the pressure in the centre is $3\,\%$ below Hertz on both meshes
(modelling error), while the contact radius and the pressure peak at the edge of the contact zone
improve with the resolution: on the coarse mesh, the active set boundary cuts only four elements
across $a$ and a local pressure overshoot of factor two appears at the edge, which vanishes after
one refinement.

\begin{table}[ht]
  \centering
  \begin{tabular}{lrrrrr}
    \toprule
    mesh & DoFs & $a/h$ & err.\ $a$ & err.\ $p(0)$ & $\max\lambda / p_0$ \\
    \midrule
    $L=6$ & 41\,865 & 4 & $-5.1\,\%$ & $-2.9\,\%$ & 2.10 \\
    $L=7$ & 72\,315 & 8 & $-1.3\,\%$ & $-3.0\,\%$ & 1.03 \\
    \bottomrule
  \end{tabular}
  \caption{Hertz problem in 3d (sphere), $\delta = 10^{-3}$, $\varepsilon = 10^5$.}
  \label{tab:hertz3d}
\end{table}

\section{Nonlinear and linear solvers}

\begin{figure}[ht]
  \centering
  \begin{tikzpicture}
    \begin{axis}[width=0.48\textwidth, height=0.4\textwidth, ymode=log, grid=major,
        grid style={gray!25}, xlabel={Newton iteration}, ylabel={$\|r\|$},
        legend style={font=\footnotesize}, legend pos=south west,
        tick label style={font=\footnotesize}, label style={font=\small}]
      \addplot+[mark=*] table[x=it, y=load07] {data/newton.dat};
      \addlegendentry{load factor 0.7}
      \addplot+[mark=square*] table[x=it, y=load10] {data/newton.dat};
      \addlegendentry{load factor 1.0}
    \end{axis}
  \end{tikzpicture}
  \caption{Newton residual (2d Hertz, $\delta=10^{-2}$, $\varepsilon=10^4$). After one or two
  iterations in which the active set changes (one of them damped), the convergence is quadratic,
  which confirms the consistency of the linearisation including the curvature term.}
  \label{fig:newton}
\end{figure}

\Cref{fig:newton} shows the Newton residual for two load steps of the Hertz problem. Once the
active set is fixed, the residual decreases quadratically (e.g.\
$9.4\cdot10^{-6}\to1.9\cdot10^{-8}\to1.1\cdot10^{-13}$). The total number of Newton iterations
grows with the penalty, from 30 for $\varepsilon=10^2$ to 175 for $\varepsilon=10^7$
(\cref{tab:penalty}), since the active set reacts more strongly to small displacement changes for
stiff penalties.

\begin{table}[ht]
  \centering
  \pgfplotstableread{data/penalty.dat}\penaltytab
  \pgfplotstabletypeset[
    columns={eps,penetration_over_delta,newton,max_linear},
    columns/eps/.style={column name={$\varepsilon$}, sci, sci zerofill, precision=0},
    columns/penetration_over_delta/.style={column name={$\max|g|/\delta$}, sci, precision=1},
    columns/newton/.style={column name={Newton it.\ (total)}},
    columns/max_linear/.style={column name={max.\ avg.\ CG it.}, fixed, precision=1},
    every head row/.style={before row=\toprule, after row=\midrule},
    every last row/.style={after row=\bottomrule},
  ]{\penaltytab}
  \caption{Penalty study (2d Hertz, $\delta=10^{-3}$, AMG): penetration and iteration counts.}
  \label{tab:penalty}
\end{table}

With AMG (matrix-based) the linear iteration counts stay small. With the matrix-free $hp$-multigrid
in 3d and $\varepsilon=10^5$, however, the average number of CG iterations grows from 5 in the
first load step (no contact) to about 130 when the contact zone is fully developed. The penalty
term adds a stiffness $\varepsilon\gg E/h$ on the contact faces, which the Chebyshev--Jacobi
smoother and in particular the coarse levels (on which the contact zone is under-resolved)
represent poorly.

\section{Summary and next steps}

\begin{itemize}
  \item The penalty contact with rigid obstacles is consistent: exact patch test, convergence in
    $h$ and $\varepsilon^{-1}$, quadratic Newton convergence; the deviation from Hertz is the
    expected finite-strain effect and vanishes like $\delta^{1/2}$.
  \item The method depends on the choice of $\varepsilon$: too small leads to visible
    penetration, too large to pressure oscillations at the quadrature points, more Newton
    iterations ($30\to175$ for $\varepsilon=10^2\to10^7$) and growing multigrid iterations.
  \item Next step: Nitsche's method (consistent, without a user-tuned penalty error and with a
    penalty of size $\gamma E/h$ only), which should also improve the multigrid performance.
\end{itemize}

\end{document}
